% VI36-DETAILED-PROBLEM-03
\begin{problem}[VI.36.P3: Affine chart and hyperplane at infinity]\label{prob:vi36-03}
Show that choosing any nonzero linear form \(L\) on \(k^{n+1}\) produces an affine open \(D_+(L)\cong\mathbb A^n_k\) whose complement is a projective hyperplane.
\end{problem}

\ifdefined\FullProblemDossiers
\begin{solution}
A nonzero linear form
\[
L=a_0x_0+\cdots+a_nx_n
\]
can be completed to a basis of the dual vector space. Therefore there is an invertible linear change of variables
sending \(L\) to \(x_0\). The corresponding matrix lies in \(\operatorname{GL}_{n+1}(k)\), hence induces a graded
automorphism of
\[
k[x_0,\ldots,x_n]
\]
and therefore an automorphism of \(\mathbb P^n_k\).

Under this automorphism,
\[
D_+(L)
\]
is carried to
\[
D_+(x_0)\cong\mathbb A^n_k.
\]
Its complement is carried to
\[
V_+(x_0),
\]
which is the projective linear locus with coordinates
\[
[0:x_1:\cdots:x_n]
\]
and therefore isomorphic to \(\mathbb P^{n-1}_k\). Thus
\[
\boxed{
D_+(L)\cong\mathbb A^n_k,
\qquad
V_+(L)\cong\mathbb P^{n-1}_k.
}
\]
\end{solution}
\fi
